\documentclass[10pt,a4paper]{amsart}
\usepackage[margin=19mm]{geometry}
\usepackage{amsmath,amssymb,mathtools}
\usepackage[scr=rsfs,cal=euler]{mathalfa}
\usepackage{xfrac}
\usepackage[unicode,hidelinks,bookmarks=false]{hyperref}
\newcommand{\ie}{i.e.\ }
\newcommand{\cf}{cf.\ }
\newcommand{\resp}{respectively\ }
\newcommand{\Subsection}{\subsection}
\providecommand{\nobreakdash}{\nobreak\hbox{-}}
\setlength{\emergencystretch}{2em}
\multlinegap0pt
\usepackage{amssymb}
\newtheorem{asA}{A}
\newtheorem{lem}{Lemma}
\def\R{\mathbb{R}}
\def\d{\delta}
\def\df{\mathrm{d}}
\def\l{\left}
\def\ol#1{\overline{#1}}
\def\p{\partial}
\def\r{\right}
\def\th{\theta}
\def\toP{\stackrel{p}{\longrightarrow}}
\def\wh#1{\widehat{#1}}
\title{M-estimation for Gaussian processes with time-inhomogeneous drifts from high-frequency data}
\author{Yasutaka SHIMIZU}
\date{Source: arXiv:2508.01164v4 (CC BY 4.0)}
\begin{document}
\maketitle

\noindent\emph{Source and licence.} M-estimation for Gaussian processes with time-inhomogeneous drifts from high-frequency data. Version arXiv:2508.01164v4 (CC BY 4.0), distributed by arXiv under the Creative Commons Attribution 4.0 International licence, http://creativecommons.org/licenses/by/4.0/, as recorded in the arXiv licence metadata for this exact version and shown on the abstract page https://arxiv.org/abs/2508.01164v4. Full licence notice: \texttt{licenses/arxiv-2508.01164-CC-BY-4.0.txt}.\par\smallskip

\noindent\textit{Editorial scope.} Counted unit: the article's lem lem:GY (source0.tex lines 2407--2416). The article's complete original proof of it is delivered with it (source0.tex lines 2418--2486, one \texttt{proof} environment of 69 lines). No mathematics is written, repaired or re-derived: the statement and the proof are the article's own lines, in the article's own notation, and no retained line is substituted or re-flowed.  Retained uncounted as the source-local context the proof needs: lines 287--289; lines 291--296; lines 1839--1851.  Nothing was omitted from any retained span, so the package carries no omission record.  No retained line contains a citation, so the wrapper carries no bibliography and no bibliography entry.  No retained line contains a figure environment, an image-inclusion command or drawing code, so no image is delivered and the package declares no asset dependency.  Editorial formatting only, in addition: an AMS \texttt{amsart} preamble that restates the article's own declarations and macros used by the retained lines, namely the environments \texttt{asA}, \texttt{lem} on separate counters, together with the standard packages those lines need.  The retained environments print in this document as Asa 1, Lemma 1 in order of appearance, whereas the article prints the same items with the numbering of its own full document.  The complete native source of the article as distributed in the arXiv e-print for this exact version is bundled unchanged as \texttt{sources/182603/source0.tex}.
\begin{asA}\label{as:2}
$K(t) \to 0$ as $t \to \infty$.
\end{asA}

\begin{asA}\label{as:3}
 $K \in C^2([0,\infty))$ with $\p_t K(0) = 0$ and $|\p_t^2 K(0)| > 0$, that is,
\[
K(t) = K(0) + \frac{1}{2} \p_t^2 K(0) t^2 + o(t^3), \quad t \to 0.
\]
\end{asA}

\begin{lem}\label{lem:ergod}
Let $Z = (Z_t)_{t \in \mathbb{R}}$ be a centered stationary Gaussian process with the conditions A\ref{as:2} and  A\ref{as:3}, and 
let $f(x,\th): \R\times \ol{\Theta} \to \R$ be a measureble function such that 
\begin{align}
\sup_{\th \in \ol{\Theta}}|\p_x^k \p_\th^l f(x,\th)| \lesssim 1 + |x|^C, \label{cond:ergod}
\end{align}
for integers $k$ and $l$ with $0\le k+l\le 1$. 
Then, under the conditions that $h_n \to 0$ and $nh_n\to \infty$ as $n\to \infty$, the following hold true:
\[
\sup_{\th \in \ol{\Theta}}\l|\frac{1}{n} \sum_{i=1}^n f(Z_{t_{i-1}^n},\th) - \int_\R f(z,\th)\phi_{K(0)}(z)\,\df z\r|\toP 0,
\]
as $n\to \infty$. 
\end{lem}

\begin{lem}\label{lem:GY}
Let $G(x,y,\th) :\R^2\times \Theta\to \R$ be a function such that $G(x,\cdot,\th)\in C^2(\R)$ for each $x\in \R$ and $\th\in \Theta$.
Suppose that, for each $x\in \R$ and $k=1,2,\dots, l \ (l\ge 2)$, $\p_y^kG(x,y,\th)$ is of polynomial growth w.r.t. $y$ uniformly in $\th\in \Theta$, $G(x,x,\th) = 0$, and that  there exists a constant $M>0$ such that $|\p_y^{l+1}G(x,y,\th)| \le M$ for all $x,y\in \R$ and $\th \in \Theta$. Then it holds that 
\begin{align*}
\frac{1}{nh_n^2}\sum_{i=1}^n G(Y^n_{i-1}, Y_i^n,\th)
&\toP  \frac{\p^2 K(0)}{2 K(0)}  \int_\R \l[\p_y G(z,z)\,z - \p_y^2 G(z, z)K(0)\r] \, \phi_{K(0)}(z)\,\df z,
\end{align*}
uniformly in $\th \in \ol{\Theta}$ under $h_n\to 0$ as $n\to \infty$, where $Y_i^n= X_{t_i} - \int_0^{t_i} \mu_{\wh{\xi}_n}(s)\,\df s$ and 
$\wh{\xi}_n$ is a consistent estimator for $\xi_0$: $\wh{\xi}_n\toP \xi_0$. 
\end{lem}

\begin{proof}
Let us fix \( \theta = (\xi, \sigma) \in \ol{\Theta} \). Define
\[
T_n(\theta) := \frac{1}{n h_n^2} \sum_{i=1}^n G\l(Y_{i-1}^n, Y_i^n, \theta\r),
\]
where \( Y_i^n := X_{t_i} - \int_0^{t_i} \mu_{\wh{\xi}_n}(s)\,\df s \) with a consistent estimator \( \wh{\xi}_n \toP \xi_0 \).

By the model definition \( X_t = Z_t + \int_0^t \mu_{\xi_0}(s)\,\df s \), we have:
\[
X_{t_i} - X_{t_{i-1}} = Z_{t_i} - Z_{t_{i-1}} + \int_{t_{i-1}}^{t_i} \mu_{\xi_0}(s)\,\df s.
\]
Thus,
\[
Y_i^n - Y_{i-1}^n = (Z_{t_i} - Z_{t_{i-1}}) + r_i^n, \quad\text{where } r_i^n := \int_{t_{i-1}}^{t_i} (\mu_{\xi_0}(s) - \mu_{\wh{\xi}_n}(s))\,\df s.
\]
Using Taylor expansion of \( y \mapsto G(x,y,\theta) \) around \( y = x \), we get
\begin{align*}
G(x,y,\theta) &= \p_y G(x,x,\theta)(y - x) + \frac{1}{2} \p_y^2 G(x,x,\theta)(y - x)^2 + R(x,y,\theta),
\end{align*}
with the remainder \( R(x,y,\theta) = \frac{1}{6} \p_y^3 G(x, \zeta, \theta)(y - x)^3 \) for some \( \zeta \) between \( x \) and \( y \). Then,
\begin{align*}
T_n(\theta) &= \frac{1}{n h_n^2} \sum_{i=1}^n \Bigg\{ \p_y G(Y_{i-1}^n, Y_{i-1}^n, \theta)(Y_i^n - Y_{i-1}^n) + \frac{1}{2} \p_y^2 G(Y_{i-1}^n, Y_{i-1}^n, \theta)(Y_i^n - Y_{i-1}^n)^2 + R_i^n(\theta) \Bigg\} \\
&=: A_n(\theta)  + B_n(\theta) + C_n(\theta). 
\end{align*}
with \( R_i^n(\theta) \) as above. We now evaluate \( A_n(\theta) \) in detail:
\begin{align*}
A_n(\theta) &= \frac{1}{n h_n^2} \sum_{i=1}^n \p_y G(Y_{i-1}^n, Y_{i-1}^n, \theta)(\Delta_i^n Z + r_i^n) \\
&= \frac{1}{n h_n^2} \sum_{i=1}^n \p_y G(Z_{t_{i-1}}, Z_{t_{i-1}}, \theta) \Delta_i^n Z + R'_n(\theta),
\end{align*}
where the error term \( R'_n(\theta) \) is decomposed as
\begin{align*}
R'_n(\theta) &= \frac{1}{n h_n^2} \sum_{i=1}^n \l[ \p_y G(Y_{i-1}^n, Y_{i-1}^n, \theta) - \p_y G(Z_{t_{i-1}}, Z_{t_{i-1}}, \theta) \r] \Delta_i^n Z \\
&\quad + \frac{1}{n h_n^2} \sum_{i=1}^n \p_y G(Y_{i-1}^n, Y_{i-1}^n, \theta) r_i^n.
\end{align*}
For the first term, using the integral form of the mean value theorem, we write:
\begin{align*}
\p_y G(Y_{i-1}^n, Y_{i-1}^n, \theta) - \p_y G(Z_{t_{i-1}}, Z_{t_{i-1}}, \theta) = \int_0^1 \nabla_1 \p_y G(Z_{t_{i-1}} + u \delta_{i-1}^n, Z_{t_{i-1}} + u \delta_{i-1}^n, \theta) \cdot \delta_{i-1}^n \,\df u,
\end{align*}
where $\d_i^n = \int_0^{t_i} [\mu_{\xi_0}(s) - \mu_{\wh{\xi}_n}(s)]\,\df s$, so that the difference is bounded as
\[
\l| \p_y G(Y_{i-1}^n, Y_{i-1}^n, \theta) - \p_y G(Z_{t_{i-1}}, Z_{t_{i-1}}, \theta) \r| \le C(1 + |Z_{t_{i-1}}|^k) |\delta_{i-1}^n|,
\]
for some constant \( C \) and \( k \ge 0 \), using the polynomial growth of \( \nabla_1 \p_y G \).
By Lemma~\ref{lem:ergod}, this yields
\[
\frac{1}{n h_n^2} \sum_{i=1}^n \l| \p_y G(Y_{i-1}^n, Y_{i-1}^n, \theta) - \p_y G(Z_{t_{i-1}}, Z_{t_{i-1}}, \theta) \r| |\Delta_i^n Z| = O_p(h_n) \to 0.
\]
The second term is bounded by
\[
\l| \frac{1}{n h_n^2} \sum_{i=1}^n \p_y G(Y_{i-1}^n, Y_{i-1}^n, \theta) r_i^n \r| \le \frac{C}{n h_n} \sum_{i=1}^n (1 + |Z_{t_{i-1}}|^k) o_p(h_n)= o_p(1),
\]
again using Lemma~\ref{lem:ergod} and the consistency \( \wh{\xi}_n \toP \xi_0 \).
Therefore,  using Lemma \ref{lem:ergod} again, we conclude that 
\[
A_n(\theta) \toP -\frac{\p^2 K(0)}{2K(0)} \int_{\R} \p_y G(z, z, \theta) z \phi_{K(0)}(z)\,\df z, 
\]
The term $B_n(\theta)$ can be evaluated similarly using ergodic theory and standard limit theorems.

As for $C_n(\theta)$, recall that
\[
R_i^n(\theta) = \frac{1}{6} \p_y^3 G\l(X_{t_{i-1}}, \zeta_i, \theta \r) (X_{t_i} - X_{t_{i-1}})^3.
\]
By the assumption that $|\p_y^3 G(x,y,\theta)| \le M$ and the fact that $X_{t_i} - X_{t_{i-1}} = O_p(h_n^{1/2})$, we get
\begin{align*}
|C_n(\theta)| &\le \frac{M}{6 n h_n^2} \sum_{i=1}^n |X_{t_i} - X_{t_{i-1}}|^3 \\
&= O_p\l(\frac{1}{n h_n^2} \sum_{i=1}^n h_n^{3/2} \r) = O_p(h_n^{1/2}) \to 0.
\end{align*}
Hence, the result follows. 
\end{proof}

\end{document}
